\chapter{Maximum-Likelihood Estimation}
\marginpar{\small\textbf{Feb. 10}}

Finding an UMVUE can be mathematically challenging, and in some cases, it might not even exist. More importantly, while an UMVUE minimizes variance
among unbiased estimators, strictly adhering to the unbiased constraint can somestimes lead to a significantly higher variance than a slightly biased estimator.
\subsubsection{Illuminating examples}
Perhaps the best known example
\[
    X \sim Poisson(\lambda),
\]
where $\lambda > 0$ is the unknown parameter. If our goal is to estimate
\[
    g(\theta) = \exp(-2\lambda),
\]
then it can be shown that the only unbiased estimator is give by the laughable
\[
    \hat{g}(X) = (-1)^X
\]
Another illuminating example is to estimate the location of a light souce based on its projection on the boundary of
a dart board over a randomly placed dart.
\section{Maximum-likelihood estimator (MLE)}
An alternative frequentist estimator is the maximum-likelihood estimator (MLE). Focusing on estimating the model parameter $\theta$ itself. Given a likelihood
family $f_\theta(x)$, an estimator $\hat{\theta} = \hat{\theta}(X)$ is an MLE of $\theta$ if
\[
    \hat{\theta}(x) \in \arg\max_{\theta\in\Theta} f_\theta(x), \quad \forall x\in\Xc
\]
or, equivalently,
\[
    \hat{\theta}(x) \in \arg\max_{\theta\in\Theta} \log f_\theta(x), \quad \forall x\in\Xc
\]
\subsection{Maximizing likelihood and minimizing K-L divergence}
While mainly motivated by intuition, MLE has an interesting interpretation when the observation $X = (X_1, \dots, X_n)$, where
\[
    X_i \iid p_\theta(x_i)
\]
More specifically, given an observation $x = (x_1, \dots, x_n)$, let
\[
    \hat{p}_x (z) = \frac{1}{n}\sum_{i=1}^n 1_{\{x_i=z\}}, \quad \forall z \in\Xc
\]
be the \textit{empirical} estimate of the \textit{marginal} distribution given the observation $x$. For two pmf's $p$ and $q$
over the same alphabet $\Xc$, the \textbf{K-L divergence} between $p$ and $q$ is defined as:
\[
    D_{KL}(p || q) := \E_p \left[ \log\frac{p(Z)}{q(Z)} \right] = \sum_{z\in\Xc} p(z) \log \frac{p(z)}{q(z)}
\]
\subsubsection{Information inequality}
Occurs when
\[
    D_{KL}(p || q) \geq 0
\]
where the equality holds if and only if $p=q$. In statistics, information theory, and machine learning, one often use K-L divergence as a distance measure between two probability
distributions, even though this distance measure is generally asymmetric about the two input distributions. \\

\noindent Claim: maximizing the likelihood
\[
    \sum_{i=1}^n \log p_\theta(x_i)
\]
is equivalent to minimzing the K-L divergence
\[
    D_{KL}(\hat{p}_x || p_\theta)
\]
\begin{proof}
    First note that
    \begin{align*}
        D_{KL}(\hat{p}_x || p_\theta) & = \E_{\hat{p}_x} \left[ \log \frac{\hat{p}_x(Z)}{p_\theta(Z)}\right]   \\
                                      & = \E_{\hat{p}_x} [\log\hat{p}_x(Z)] - \E_{\hat{p}_x}[\log p_\theta(Z)]
    \end{align*}
    Further note that
    \begin{align*}
        \E_{\hat{p}_x} & = \sum_{z\in\Xc} \hat{p}_x(z)\log p_{\theta}(z)                               \\
                       & = \frac{1}{n} \sum_{z\in\Xc}\sum_{i=1}^n 1_{\{ x_i=z \}} \log p_\theta(z)     \\
                       & = \frac{1}{n} \sum_{i=1}^n \bp{\sum_{z\in\Xc} 1_{\{x_i=z\}} \log p_\theta(z)} \\
                       & = \frac{1}{n} \sum_{i=1}^n \log p_\theta(x_i)
    \end{align*}
    Combining the above two results gives
    \[
        D_{KL}(\hat{p}_x || p_\theta) = \E_{\hat{p}_x}[\log \hat{p}_x (Z)] - \frac{1}{n} \sum_{i=1}^n \log p_\theta(x_i)
    \]
    We may thus conclude that minimizing the K-L divergence $D_{KL}(\hat{p}_x || p_\theta)$ is equivalent to
    maximizing the likelihood $\sum_{i=1}^n \log p_\theta(x_i)$. The above result also suggests that if
    our model is identifiable, MLE will recover $\theta$ in the limit as $n \to \infty$. We will have more
    to say about this result a bit later.
\end{proof}

\subsection{Re-parameterization}
Assume that the likelihood function $f_\theta(x)$ can be re-parameterized in terms of a more natural parameter
\[
    \eta = \eta(\theta)
\]
In this case, for any given observation $x$, we have
\[
    \arg\max_{\theta\in\Theta} f_\theta(x) = \{ \theta \in \Theta : \eta(\theta) \in \arg\max_{\theta\in\Theta} \}
\] ... %TODO
Consider an exponential family of distribution re-parameterized by the natural parameter $\eta$:
\[
    \tilde{f}_\eta(x) = h(x) \exp(\eta^t T(x) - A(\eta))
\]
where the space of the natural parameter is an open and connected subset of $\R^s$. To find an MLE for $\eta$, note that
\begin{align*}
    \nabla_\eta \log\tilde{f}_\eta(x)    & = T(x) - \nabla_\eta A(\eta) = T(x) - \E_\eta[T(X)] \\
    \nabla_\eta^2 \log \tilde{f}_\eta(x) & = -\nabla_\eta^2 A(\eta) = \Cov_\eta[T(X)]
\end{align*}
A covariance mastrix is always positive semi-definite, so for any $x\in\Xc$, the log-likelihood $\log\tilde{f}_\eta(x)$ is a concave function of $\eta$, for which any stationary point is a global maxima. Therefore,
any solution for $\eta$ to the equation
\[
    \E_\eta[T(X)] = T(x)
\]
is an MLE for $\eta$. Furthermore, if the covariance matrix of $T(X)$ is positive definite for all
$\eta$, for any $x \in \Xc$ the log-likelihood $\log\tilde{f}_\eta(x)$ is a strictly concave function of $\eta$. In this case, the MLE is unique if it exists.

Finally, by the re-parameterization property, any solution for $\theta$ to the equation
\[
    \E_\theta[T(X)] = T(x)
\]
is an MLE for $\theta$.
\subsubsection{Gaussian with unknown mean and variance}
Recall that in this case, a sufficient statistic is given by
\[
    T_1 = \sum_{i=1}^n X_i \quad\text{and}\quad T_2 = \sum_{i=1}^n X_i^2
\]
the expectations of the sufficient statistics are given by:
\[
    \E_\theta[T_1] = n\mu \quad\text{and}\quad \E_\theta[T_2] = n(\sigma^2 + \nu^2)
\]
From previous discussion, any solution of $(\mu, \sigma^2)$ to the equations
\begin{align*}
    n\mu                & = T_1(x) \\
    n(\sigma^2 + \mu^2) & = T_2(x)
\end{align*}
is an MLE. The above equations have a unique solution given by
\[
    \mu = \frac{T_1(x)}{n} \quad\text{and}\quad \sigma^2 = \frac{T_2(x)}{n} - \bp{\frac{T_1(x)}{n}}^2
\]
We may thus conclue that the empirical mean
\[
    \hat{\mu} = \frac{T_1(x)}{n} = \frac{1}{n} \sum_{i=1}^n X_i
\]
is the MLE for $\mu$ and the empirical variance
\[
    \hat{\sigma}^2 = \frac{T_2(x)}{n} - \bp{\frac{T_1(x)}{n}}^2 = \frac{1}{n}\sum_{i=1}^n X_i^2 - \hat{\mu}^2 = \frac{1}{n}\sum_{i=1}^n (X_i - \hat{\mu})^2
\]
is the MLE for $\sigma^2$, when both $\mu$ and $\sigma^2$ are unknown. By Cochran's theorem,
\[
    \E_\theta \left[ \sum_{i=1}^n (X_i - \hat{\mu})^2\right] = (n-1)\sigma^2
\]
so
\[
    \E_\theta[\hat{\sigma^2}] = \frac{n-1}{n}\sigma^2
\]
i.e., the empirical variance is a biased estimator for the true variance $\sigma^2$.

\subsection{Consistency and asymptotic normality}
Note that in the previous example, while the MLE for $\sigma^2$ is biased, the bias vanishes asymptotically in the limit as the sample size $n\to\infty$. Also note that the variances of the MLE for both $\mu$ and $\sigma^2$ coincide with the CRLB for unbiased estimators. Under certain regularity conditions, it can be shown that the above observations hold for a general likelihood function.

More specifically, let $X = (X_1, \dots, X_n)$, where $X_i \iid f_\theta(x_i)$. For any $\theta$ that is in the interior of $\Theta$, let $\hat{\theta}$ be the MLE for $\theta$ given by the unique solution to the equation
\[
    \nabla_\theta \bp{\sum_{i=1}^n \log f_\theta(x_i)} = 0
\]
We have $\sqrt{n}(\hat{\theta} - \theta)$ converges in distribution to $\Nc(0, I_1^{-1}(\theta))$ in the limit as the sample size $n\to\infty$.
That is, when the sample size $n$ is large, the MLE $\hat{\theta}$ is approximately distributed as
\[
    \Nc \bp{\theta, \frac{1}{n}I_1^{-1}(\theta)}
\]
or, equivalently,
\[
    \Nc (\theta,I_n^{-1}(\theta))
\]
due to the additivity of Fisher information $I_n(\theta) = n I_1(\theta)$. Therefore, under certain regularity conditions, the MLE is consistent in that it can recover the unknown parameter $\theta$ in the limit as the sample size $n\to\infty$. The proof of the asymptotic normality of the MLE is based on the well-known central limit theorem. This is omitted here, however.

\subsection{Multinomial with unknown probabilities}
Goal is to use the empirical counts to estimate the true probability.
Let $X\sim p_\theta(x)$, where the observation
\[
    X = (N_1, \dots, N_K)
\]
is the empirical counts over $K$ categories and $n$ total counts, the unknown parameter
\[
    \theta = (p_1, \dots, p_K)
\]
is a probability vector over $K$ categories, and the likelihood function is multinomial:
\[
    p_\theta(x) = \frac{n!}{n_1!\cdots n_K!} \prod_{k=1}^K{p_k^{n_k}}
\]
% TODO: finish this whole section 
The total counts $n$ and the number of categories $K$ are assumed to be known. The MLE can be found by
solving the equation
\[
    \nabla_\theta L_\lambda (\theta) = 0
\]
and then finding a Lagrangian multiplier $\lambda \geq 0$ to satisfy the constraint
\[
    \sum_{k=1}^K p_k = 1
\]
The MLE for $\theta$ in this case is given by the empirical frequency estimator
\[
    \hat{\theta} = \bp{\frac{N_1}{n}, \dots, \frac{N_K}{n}}
\]
Note that for any $k \in [K]$
\[
    \E_\theta[N_k] = np_k
\]
so the empirical frequency estimator $\hat{\theta}$ is unbiased.
\section{Estimating $g(\theta)$}
Let $\eta = g(\theta)$. Note that if the likelihood function $f_\theta(x)$ can be rep-parameterized in terms of
$\eta$, then it is natural to use an MLE of $\eta$ as an "MLE" estimate of $g(\theta)$. By the re-parameterization
property, $\eta^*$ is an MLE of $\eta$ if and only if
\[
    \eta^* = g(\theta^*)
\]
for some $\theta^*$ that is an MLE of $\theta$. Therefore, if the likelihood function $f_\theta(x)$ can be re-parameterized
in terms of $g(\theta)$, the "MLE" of $g(\theta)$ is given by $g(\theta^*)$, where $\theta^*$ is an MLE of $\theta$.

In statistics, it is customary to define the MLE of $g(\theta)$ as $g(\theta^*)$, whether the likelihood function $f_\theta(x)$ can be re-parameterized
in terms of $g(\theta)$ or not. This is known as the invariance property of MLE. Many other estimators
such as UMVUE do not enjoy the invariance property.